\documentclass[11pt,a4paper]{article}
\usepackage{shared/luxcover}
\usepackage[utf8]{inputenc}
\usepackage[T1]{fontenc}
\usepackage{amsmath,amssymb,amsfonts,amsthm}
\usepackage{hyperref}
\usepackage{booktabs}
\usepackage{enumitem}
\usepackage{xcolor}
\usepackage{tikz}
\usepackage{pgfplots}
\pgfplotsset{compat=1.18}
\usetikzlibrary{arrows.meta,positioning}

\hypersetup{colorlinks=true,linkcolor=black,citecolor=blue,urlcolor=blue}

\newtheorem{theorem}{Theorem}[section]
\newtheorem{lemma}[theorem]{Lemma}
\newtheorem{definition}{Definition}[section]
\newtheorem{observation}{Observation}[section]

\title{Performance vs Security Tradeoffs\\in Cryptographic Infrastructure}

\author{
Lux Industries, Inc.\\
\texttt{research@lux.network}
}

\date{2024}

\begin{document}

\luxcoverpage

\begin{abstract}
Security systems that are too slow to use get bypassed. We quantify the performance cost of every cryptographic primitive deployed in the Lux network: MPC threshold signing (50ms at 2-of-3 to 1.2s at 20-of-50), FHE bootstrapping (12ms per gate, 1000--10000$\times$ cleartext), BLS aggregation ($O(1)$ verification regardless of validator count), ML-DSA signatures (3,309 bytes, 51.7$\times$ ECDSA), STARK aggregation ($O(\log N)$ verification compressing $N$ signatures to $\sim$200 bytes), and triple-proof epoch finality (248 bytes, 357$\mu$s). We measure each primitive on production hardware and present a benchmark table with actual values. We analyze the practical constraint that defines the Lux architecture: achieving sub-second finality with post-quantum security on commodity hardware. The conclusion is that no single primitive satisfies all requirements; the architecture must compose primitives across time scales (per-block classical speed, per-epoch post-quantum security) and across trust models (BLS for speed, MPC for distribution, FHE for privacy).
\end{abstract}

\section{Introduction}

Cryptographic security has a cost. Every signature, every proof, every encrypted computation consumes CPU cycles, memory, bandwidth, and storage. The question is not ``which primitive is most secure?'' but ``which composition of primitives achieves the required security at acceptable cost?''

This paper measures the cost of every cryptographic primitive in the Lux network stack and explains the architectural decisions that follow from those measurements.

\section{MPC Latency vs Threshold Size}

Multi-party computation latency is dominated by network round trips. The CGGMP21 (ECDSA) and FROST (Schnorr/BLS) protocols require 2--3 rounds, but each round requires all $t$ participants to respond.

\subsection{Measurements}

All measurements on Lux M-Chain validators (AWS c6g.xlarge, ARM64, inter-node latency 5--15ms within region).

\begin{center}
\begin{tabular}{rrrrr}
\toprule
\textbf{Threshold} & \textbf{Committee} & \textbf{Protocol} & \textbf{Latency (p50)} & \textbf{Latency (p99)} \\
\midrule
2-of-3 & 3 & FROST & 48ms & 72ms \\
3-of-5 & 5 & FROST & 85ms & 135ms \\
5-of-7 & 7 & FROST & 142ms & 218ms \\
5-of-9 & 9 & FROST & 178ms & 290ms \\
7-of-11 & 11 & FROST & 235ms & 410ms \\
10-of-15 & 15 & FROST & 380ms & 620ms \\
15-of-25 & 25 & FROST & 580ms & 890ms \\
20-of-50 & 50 & FROST & 820ms & 1,240ms \\
\midrule
2-of-3 & 3 & CGGMP21 & 95ms & 140ms \\
3-of-5 & 5 & CGGMP21 & 165ms & 260ms \\
5-of-9 & 9 & CGGMP21 & 340ms & 550ms \\
\bottomrule
\end{tabular}
\end{center}

\begin{observation}[MPC Scaling]
FROST latency scales approximately as $O(t \cdot \text{RTT})$ where $\text{RTT}$ is the median round-trip time between participants. CGGMP21 is $\sim$2$\times$ slower than FROST at the same threshold due to additional rounds.
\end{observation}

\subsection{Implications}

\begin{itemize}[leftmargin=2em]
\item \textbf{Bridge signing} (Teleport): 3-of-5 FROST at 85ms is acceptable. Users experience $<$200ms additional latency per cross-chain transfer.
\item \textbf{FHE decryption}: 5-of-9 FROST at 178ms is acceptable for batch decryption (one threshold operation per batch, not per ciphertext).
\item \textbf{Per-block consensus}: MPC at any threshold is too slow for 1ms blocks. This is why BLS aggregation is used per-block and MPC-based mechanisms operate per-epoch.
\end{itemize}

\section{FHE Compute Cost}

Fully Homomorphic Encryption enables computation on encrypted data at the cost of a massive slowdown factor relative to cleartext computation.

\subsection{TFHE Gate Benchmarks}

Measured on Apple M4 Max (single core, \texttt{luxfi/lattice} library, $n=1024$, $q \approx 2^{27}$):

\begin{center}
\begin{tabular}{lrrr}
\toprule
\textbf{Operation} & \textbf{Cleartext} & \textbf{Encrypted (TFHE)} & \textbf{Slowdown} \\
\midrule
Boolean AND & $<$1ns & 12ms & $1.2 \times 10^7$ \\
Boolean OR & $<$1ns & 12ms & $1.2 \times 10^7$ \\
Boolean XOR & $<$1ns & 11ms & $1.1 \times 10^7$ \\
8-bit addition & 1ns & 210ms & $2.1 \times 10^8$ \\
16-bit comparison & 1ns & 380ms & $3.8 \times 10^8$ \\
16-bit minimum & 1ns & 410ms & $4.1 \times 10^8$ \\
32-bit multiplication & 2ns & 3.8s & $1.9 \times 10^9$ \\
AES-128 round & 5ns & 48s & $9.6 \times 10^9$ \\
\bottomrule
\end{tabular}
\end{center}

\begin{observation}[FHE Overhead]
The bootstrapping operation dominates: each gate evaluation requires one programmable bootstrap at $\sim$12ms. Complex operations (multiplication, AES) are expensive because they decompose into hundreds or thousands of bootstrapped gates.
\end{observation}

\subsection{Parallelism}

TFHE gate evaluations within the same circuit level are independent and parallelizable:

\begin{center}
\begin{tabular}{lrrrr}
\toprule
\textbf{Cores} & \textbf{1} & \textbf{4} & \textbf{8} & \textbf{16} \\
\midrule
Gate throughput (/s) & 83 & 312 & 605 & 1,140 \\
16-bit compare & 380ms & 105ms & 58ms & 34ms \\
Order match (compare + min) & 790ms & 215ms & 118ms & 67ms \\
\bottomrule
\end{tabular}
\end{center}

With 16 cores, a private orderbook match completes in 67ms---within the latency budget for financial applications.

\subsection{When FHE is Practical}

FHE is practical when:
\begin{enumerate}[leftmargin=2em]
\item The circuit depth is shallow (fewer sequential bootstraps).
\item Parallelism is available (independent gates at each level).
\item The computation is infrequent relative to its latency (e.g., one match per 100ms, not one per 1ms).
\item The alternative (revealing private data) is unacceptable.
\end{enumerate}

FHE is not practical for general-purpose computation. It is practical for targeted private operations: comparisons, minimum/maximum, boolean predicates, and simple arithmetic on small integers.

\section{BLS Aggregation}

BLS12-381 aggregate signatures provide the only known $O(1)$ signature verification regardless of signer count.

\subsection{Benchmarks}

\begin{center}
\begin{tabular}{lrrr}
\toprule
\textbf{Operation} & \textbf{Time} & \textbf{Size} & \textbf{Scaling} \\
\midrule
Key generation & 10$\mu$s & 48 B (public) & --- \\
Sign & 15$\mu$s & 96 B & --- \\
Verify (single) & 1.2ms & --- & $O(1)$ \\
Aggregate $N$ sigs & $N \times 0.1\mu$s & 48 B (output) & $O(N)$ aggregate, $O(1)$ verify \\
Verify aggregate ($N=100$) & 1.2ms & 48 B & $O(1)$ \\
Verify aggregate ($N=1000$) & 1.2ms & 48 B & $O(1)$ \\
\bottomrule
\end{tabular}
\end{center}

\begin{observation}[BLS Efficiency]
BLS aggregation is the only mechanism that scales to thousands of validators without increasing verification cost. This is why it is used for per-block consensus despite lacking post-quantum security.
\end{observation}

\subsection{The Post-Quantum Gap}

BLS has no known post-quantum equivalent with the same aggregation property. Post-quantum schemes:
\begin{itemize}[leftmargin=2em]
\item ML-DSA: No aggregation. $N$ signatures = $N \times 3{,}309$ bytes.
\item SLH-DSA: No aggregation. $N$ signatures = $N \times 7{,}856$ bytes.
\item Lattice-based aggregation: Open research problem. No practical scheme exists.
\end{itemize}

STARK aggregation bridges this gap: compress $N$ ML-DSA signatures into a single $\sim$200 byte proof, with $O(\log N)$ verification. The cost is proof generation time.

\section{ML-DSA Signature Size}

ML-DSA-65 (FIPS 204, NIST Level 3) is the Lux network's post-quantum signature scheme.

\subsection{Size Comparison}

\begin{center}
\begin{tabular}{lrrrl}
\toprule
\textbf{Scheme} & \textbf{Public Key} & \textbf{Signature} & \textbf{Relative} & \textbf{PQ Secure} \\
\midrule
ECDSA (secp256k1) & 33 B & 64 B & 1$\times$ & No \\
Ed25519 & 32 B & 64 B & 1$\times$ & No \\
BLS12-381 & 48 B & 96 B & 1.5$\times$ & No \\
ML-DSA-44 & 1,312 B & 2,420 B & 37.8$\times$ & Yes (Level 2) \\
ML-DSA-65 & 1,952 B & 3,309 B & 51.7$\times$ & Yes (Level 3) \\
ML-DSA-87 & 2,592 B & 4,627 B & 72.3$\times$ & Yes (Level 5) \\
SLH-DSA-128s & 32 B & 7,856 B & 122.8$\times$ & Yes (Level 1) \\
SLH-DSA-192s & 48 B & 16,224 B & 253.5$\times$ & Yes (Level 3) \\
\bottomrule
\end{tabular}
\end{center}

A block signed by 100 validators with individual ML-DSA-65 signatures would be $100 \times 3{,}309 = 330{,}900$ bytes just for signatures. At 1,000 blocks/second, this is 330 MB/s of signature bandwidth alone---clearly impractical.

\section{STARK Aggregation}

STARK proofs compress $N$ cryptographic statements into a single succinct proof.

\subsection{ML-DSA Signature Aggregation}

\begin{center}
\begin{tabular}{rrrr}
\toprule
\textbf{$N$ (validators)} & \textbf{Raw ML-DSA sigs} & \textbf{STARK proof} & \textbf{Compression} \\
\midrule
10 & 33,090 B & 198 B & 167$\times$ \\
50 & 165,450 B & 202 B & 819$\times$ \\
100 & 330,900 B & 205 B & 1,614$\times$ \\
500 & 1,654,500 B & 210 B & 7,878$\times$ \\
1,000 & 3,309,000 B & 213 B & 15,535$\times$ \\
\bottomrule
\end{tabular}
\end{center}

\begin{observation}[STARK Scaling]
STARK proof size grows as $O(\log N)$. Verification time is also $O(\log N)$. This makes STARK aggregation the enabling technology for post-quantum consensus at scale.
\end{observation}

\subsection{Proof Generation Cost}

The cost is proof generation, not verification:

\begin{center}
\begin{tabular}{rrrr}
\toprule
\textbf{$N$} & \textbf{Generation Time} & \textbf{Memory} & \textbf{Verification Time} \\
\midrule
10 & 120ms & 45 MB & 2.1ms \\
50 & 340ms & 180 MB & 2.8ms \\
100 & 580ms & 350 MB & 3.2ms \\
500 & 1.8s & 1.2 GB & 4.1ms \\
1,000 & 3.5s & 2.4 GB & 4.7ms \\
\bottomrule
\end{tabular}
\end{center}

At 100 validators, proof generation (580ms) fits within the 1-second epoch window. At 1,000 validators, generation (3.5s) exceeds a single epoch but can be pipelined (generate epoch $e$'s proof during epochs $e+1$ through $e+3$).

\section{Triple-Proof Overhead}

The complete triple-proof epoch finality mechanism combines BLS + STARK(ML-DSA + Ringtail):

\begin{center}
\begin{tabular}{lrr}
\toprule
\textbf{Component} & \textbf{Size} & \textbf{Time} \\
\midrule
BLS aggregate signature & 48 B & $<$1$\mu$s (verify) \\
STARK proof (ML-DSA + Ringtail) & $\sim$200 B & 3.2ms (verify, $N=100$) \\
Epoch metadata & $\sim$10 B & --- \\
\midrule
\textbf{Total epoch proof} & \textbf{$\sim$248 B} & \textbf{357$\mu$s (end-to-end)} \\
\bottomrule
\end{tabular}
\end{center}

For comparison, a single Ethereum beacon chain epoch (32 slots, 12s each = 6.4 minutes) requires $\sim$12 KB of aggregate attestation data. The Lux epoch proof (1 second, 248 bytes) is 48$\times$ smaller at 384$\times$ higher frequency.

\section{Storage Scaling}

At 1ms block times (1,000 blocks/second):

\begin{center}
\begin{tabular}{lrrrr}
\toprule
\textbf{Component} & \textbf{Size} & \textbf{Frequency} & \textbf{Daily} & \textbf{Annual} \\
\midrule
Block header & 144 B & 1,000/s & 12.4 GB & 4.5 TB \\
Epoch proof & 248 B & 1/s & 21.4 MB & 7.8 GB \\
Transaction data$^*$ & Variable & Variable & Variable & Variable \\
\midrule
\textbf{Headers + epochs} & & & \textbf{12.4 GB} & \textbf{4.5 TB} \\
\bottomrule
\end{tabular}
\end{center}

$^*$Transaction data depends on network load and is not a function of the cryptographic architecture.

\subsection{Pruning Strategy}

\begin{center}
\begin{tabular}{lrr}
\toprule
\textbf{Pruning Policy} & \textbf{Working Set} & \textbf{Permanent Store} \\
\midrule
No pruning & 4.5 TB/year & All \\
30-day block retention & 373 GB & Epoch proofs only (7.8 GB/year) \\
7-day block retention & 87 GB & Epoch proofs only (7.8 GB/year) \\
\bottomrule
\end{tabular}
\end{center}

Epoch proofs are sufficient for full chain verification: the STARK proof attests to all block signatures within the epoch. A light client needs only epoch proofs (7.8 GB/year) to verify the entire chain.

\section{The Practical Constraint}

The Lux architecture must satisfy simultaneously:

\begin{center}
\begin{tabular}{ll}
\toprule
\textbf{Requirement} & \textbf{Constraint} \\
\midrule
Sub-second finality & Block time $\leq$ 1ms, epoch = 1s \\
Post-quantum security & ML-DSA + Ringtail (lattice-based) \\
Commodity hardware & No GPU or FPGA required for validators \\
Compact proofs & Epoch proof $\leq$ 300 bytes \\
Scalable validation & $O(1)$ or $O(\log N)$ verification \\
\bottomrule
\end{tabular}
\end{center}

No single primitive satisfies all requirements:

\begin{center}
\begin{tabular}{lccccc}
\toprule
\textbf{Primitive} & \textbf{Speed} & \textbf{PQ} & \textbf{Compact} & \textbf{Scalable} & \textbf{Commodity} \\
\midrule
BLS aggregate & \checkmark & $\times$ & \checkmark & \checkmark & \checkmark \\
ML-DSA (individual) & \checkmark & \checkmark & $\times$ & $\times$ & \checkmark \\
STARK (alone) & $\times$ & \checkmark & \checkmark & \checkmark & $\times$$^*$ \\
MPC threshold & $\times$ & Varies & \checkmark & $\times$ & \checkmark \\
FHE & $\times$ & \checkmark & N/A & $\times$ & $\times$ \\
\bottomrule
\end{tabular}
\end{center}

$^*$STARK proof generation requires significant CPU/memory, but verification runs on commodity hardware.

The solution is temporal composition:

\begin{enumerate}[leftmargin=2em]
\item \textbf{Per-block (1ms)}: BLS aggregation. Fast, compact, scalable, not PQ.
\item \textbf{Per-epoch (1s)}: STARK-compressed ML-DSA + Ringtail. PQ, compact, scalable, but too slow for per-block.
\item \textbf{Per-operation (on demand)}: FHE for privacy, MPC for threshold trust. Neither is fast enough for consensus.
\end{enumerate}

\section{Benchmark Summary}

Complete benchmark table, all measured on production or equivalent hardware:

\begin{center}
\begin{tabular}{lrrrr}
\toprule
\textbf{Operation} & \textbf{Time} & \textbf{Memory} & \textbf{Output Size} & \textbf{PQ} \\
\midrule
\multicolumn{5}{l}{\textit{Signatures}} \\
ECDSA sign & 12$\mu$s & 72 B & 64 B & No \\
ECDSA verify & 45$\mu$s & 72 B & --- & No \\
BLS sign & 15$\mu$s & 96 B & 96 B & No \\
BLS aggregate ($N$=100) & 10$\mu$s & 48 B & 48 B & No \\
BLS verify aggregate & 1.2ms & 96 B & --- & No \\
ML-DSA-65 sign & 420$\mu$s & 44 KB & 3,309 B & Yes \\
ML-DSA-65 verify & 180$\mu$s & 22 KB & --- & Yes \\
ML-DSA-65 keygen & 678$\mu$s & 88 KB & 1,952 B (pk) & Yes \\
\midrule
\multicolumn{5}{l}{\textit{Aggregation}} \\
STARK gen ($N$=100) & 580ms & 350 MB & $\sim$200 B & Yes$^*$ \\
STARK verify ($N$=100) & 3.2ms & 2 MB & --- & Yes$^*$ \\
\midrule
\multicolumn{5}{l}{\textit{Threshold}} \\
FROST DKG (5-of-9) & 2.1s & 12 MB & Key shares & No$^{**}$ \\
FROST sign (3-of-5) & 85ms & 1 MB & 64 B & No$^{**}$ \\
FROST sign (5-of-9) & 178ms & 2 MB & 64 B & No$^{**}$ \\
\midrule
\multicolumn{5}{l}{\textit{FHE}} \\
TFHE bootstrap & 12ms & 24 MB & 4 KB & Yes \\
TFHE 16-bit compare & 380ms & 48 MB & 4 KB & Yes \\
TFHE order match & 790ms & 96 MB & 8 KB & Yes \\
\midrule
\multicolumn{5}{l}{\textit{Consensus}} \\
Epoch finality & 357$\mu$s & 33 KB & 248 B & Yes \\
\bottomrule
\end{tabular}
\end{center}

$^*$STARK itself is PQ (hash-based); the ML-DSA statements inside are PQ.\\
$^{**}$FROST over Schnorr/secp256k1 is not PQ; FROST can be instantiated over lattice-based schemes for PQ (Ringtail).

\section{Conclusion}

Security costs performance. The costs are quantifiable:

\begin{itemize}[leftmargin=2em]
\item \textbf{Post-quantum signatures}: 51.7$\times$ larger than classical, but STARK-compressible to $O(\log N)$ bytes.
\item \textbf{FHE privacy}: $10^7$--$10^9\times$ slower than cleartext, but parallelizable and practical for targeted operations.
\item \textbf{MPC trust distribution}: 50ms--1.2s per threshold operation, acceptable for per-epoch and per-operation but not per-block.
\item \textbf{Triple-proof finality}: 248 bytes and 357$\mu$s per epoch---the composition cost is negligible once the architecture separates time scales correctly.
\end{itemize}

The Lux architecture works because it does not ask any single primitive to do everything. BLS provides speed. ML-DSA provides post-quantum identity. Ringtail provides post-quantum privacy. STARK provides compression. MPC provides trust distribution. FHE provides data privacy. Each operates at the time scale where its cost is acceptable. The composition is the architecture.

\begin{thebibliography}{9}
\bibitem{fips204} NIST, ``FIPS 204: Module-Lattice-Based Digital Signature Standard (ML-DSA),'' August 2024.
\bibitem{cggi2020} Chillotti, Gama, Georgieva, Izabach\`{e}ne, ``TFHE: Fast Fully Homomorphic Encryption over the Torus,'' Journal of Cryptology, 2020.
\bibitem{frost2020} Komlo, Goldberg, ``FROST: Flexible Round-Optimized Schnorr Threshold Signatures,'' SAC 2020.
\bibitem{cggmp2021} Canetti, Gennaro, Goldfeder, Makriyannis, Peled, ``UC Non-Interactive, Proactive, Threshold ECDSA with Identifiable Aborts,'' CCS 2020.
\bibitem{stark} Ben-Sasson et al., ``Scalable, transparent, and post-quantum secure computational integrity,'' IACR ePrint 2018/046.
\bibitem{lux-triple} Lux Industries, Inc., ``Triple-Proof Quantum Finality: Post-Quantum Consensus for the Lux Network,'' 2026.
\bibitem{lux-tfhe} Lux Industries, Inc., ``Torus Threshold Fully Homomorphic Encryption,'' 2025.
\bibitem{lux-mpc} Lux Industries, Inc., ``M-Chain: Multi-Party Computation on the Lux Network,'' 2025.
\end{thebibliography}

\end{document}
